\documentclass[11pt,a4paper]{article}
\usepackage[utf8]{inputenc}
\usepackage[T1]{fontenc}
\usepackage[french]{babel}
\usepackage{amsmath,amssymb}
\usepackage{geometry}
\usepackage{enumitem}
\usepackage{booktabs}
\usepackage{array}
\usepackage{tabularx}
\usepackage[table]{xcolor}
\usepackage{tcolorbox}
\usepackage{tikz}
\usepackage{fontawesome5}

\tcbuselibrary{skins,breakable}

\geometry{margin=1.8cm}

% Couleurs
\definecolor{primary}{RGB}{31,78,121}
\definecolor{soft}{RGB}{230,239,247}
\definecolor{amber}{RGB}{178,91,0}
\definecolor{amberbg}{RGB}{255,244,224}
\definecolor{amberborder}{RGB}{240,192,112}
\definecolor{greenbg}{RGB}{233,246,234}
\definecolor{greenborder}{RGB}{182,224,184}
\definecolor{redbg}{RGB}{253,236,236}
\definecolor{redborder}{RGB}{242,179,171}
\definecolor{green}{RGB}{46,125,50}
\definecolor{red}{RGB}{211,47,47}

% Boîtes
\newtcolorbox{idee}[1][]{
  colback=soft, colframe=primary, boxrule=1pt, arc=6pt,
  fonttitle=\bfseries\large, coltitle=white, colbacktitle=primary,
  left=8pt, right=8pt, top=6pt, bottom=6pt,
  breakable, #1
}

\newtcolorbox{retenir}[1][]{
  colback=greenbg, colframe=greenborder, boxrule=1pt, arc=6pt,
  fonttitle=\bfseries\large, coltitle=black, colbacktitle=greenborder,
  left=8pt, right=8pt, top=6pt, bottom=6pt,
  breakable, #1
}

\newtcolorbox{attention}[1][]{
  colback=redbg, colframe=redborder, boxrule=1pt, arc=6pt,
  fonttitle=\bfseries\large, coltitle=black, colbacktitle=redborder,
  left=8pt, right=8pt, top=6pt, bottom=6pt,
  breakable, #1
}

\newtcolorbox{labo}[1][]{
  colback=amberbg, colframe=amberborder, boxrule=1pt, arc=6pt,
  fonttitle=\bfseries\large, coltitle=black, colbacktitle=amberborder,
  left=8pt, right=8pt, top=6pt, bottom=6pt,
  breakable, #1
}

\title{\textcolor{primary}{\textbf{La carte de contrôle (Control Chart) — L'essentiel en 5 minutes}}\\
\large Fiche de synthèse pour étudiants BTS Bioanalyses en Laboratoire de Contrôle}
\author{}
\date{}

\begin{document}

\maketitle

\vspace{-0.5cm}

\begin{idee}[title={\faLightbulb\quad En une phrase}]
La carte de contrôle, c'est un \textbf{graphique qui surveille} un procédé de mesure dans le temps pour détecter \textbf{immédiatement} quand quelque chose déraille — avant que les résultats ne deviennent faux.
\end{idee}

\section*{\faBullseye\quad Pourquoi ça existe ?}

\begin{itemize}[leftmargin=1.2cm, itemsep=2pt]
    \item[\textcolor{primary}{\faCheck}] Pour \textbf{détecter une dérive} d'un procédé avant qu'il ne produise des résultats non conformes.
    \item[\textcolor{primary}{\faCheck}] Pour \textbf{distinguer} les variations normales (hasard) des variations anormales (cause spéciale).
    \item[\textcolor{primary}{\faCheck}] Pour \textbf{valider en routine} les séries d'analyses (contrôle qualité interne, CIQ).
    \item[\textcolor{primary}{\faCheck}] Pour \textbf{prouver la stabilité} d'un procédé (accréditation ISO 15189, ISO/IEC 17025).
\end{itemize}

\section*{\faKey\quad Les 3 lignes à connaître}

\begin{center}
\renewcommand{\arraystretch}{1.6}
\begin{tabularx}{\linewidth}{|>{\columncolor{soft}}p{4cm}|X|}
\hline
\textbf{\textcolor{primary}{\Large Ligne centrale (LC)}} & La \textbf{moyenne} de référence du procédé. \\ \hline
\textbf{\textcolor{primary}{\Large Limite supérieure (LSC)}} & Seuil au-dessus duquel le procédé est suspect. \\ \hline
\textbf{\textcolor{primary}{\Large Limite inférieure (LIC)}} & Seuil en dessous duquel le procédé est suspect. \\ \hline
\end{tabularx}
\end{center}

\vspace{0.2cm}

\begin{retenir}[title={\faBullseye\quad Règle de base}]
En général, les limites sont placées à \(\pm 3\sigma\) autour de la moyenne. Un point \textbf{hors des limites} signale une \textbf{cause spéciale} à investiguer.
\end{retenir}

\vspace{0.3cm}

% Schéma carte de contrôle
\begin{center}
\begin{tikzpicture}[scale=0.9]
  % Axes
  \draw[->, thick] (0,0) -- (11,0) node[right] {Temps};
  \draw[->, thick] (0,0) -- (0,5) node[above] {Valeur};

  % Limites
  \draw[red, thick, dashed] (0,4) -- (11,4) node[right, font=\small] {LSC};
  \draw[primary, thick] (0,2.5) -- (11,2.5) node[right, font=\small] {LC};
  \draw[red, thick, dashed] (0,1) -- (11,1) node[right, font=\small] {LIC};

  % Zones
  \fill[greenbg] (0,1) rectangle (11,4);
  \fill[soft] (0,2) rectangle (11,3);

  % Points en contrôle
  \foreach \x/\y in {0.5/2.6, 1.2/2.3, 1.9/2.7, 2.6/2.5, 3.3/2.4, 4/2.8, 4.7/2.6} {
    \fill[primary] (\x,\y) circle (0.08);
  }

  % Dérive
  \foreach \x/\y in {5.4/3.1, 6.1/3.5, 6.8/3.8, 7.5/4.3} {
    \fill[red] (\x,\y) circle (0.08);
  }

  % Point hors limite
  \fill[red] (8.2,4.6) circle (0.12);
  \node[above, red, font=\small] at (8.2,4.6) {Hors contrôle !};

  % Suite après action
  \foreach \x/\y in {9/2.6, 9.7/2.4, 10.4/2.7} {
    \fill[primary] (\x,\y) circle (0.08);
  }

  % Flèche "action corrective"
  \draw[->, thick, amber] (8.2,4.2) -- (8.8,2.8);
  \node[amber, font=\small, align=center] at (9.5,3.5) {Action\\corrective};
\end{tikzpicture}
\end{center}

\section*{\faFlask\quad Les 2 types de causes de variation}

\begin{center}
\renewcommand{\arraystretch}{1.5}
\begin{tabularx}{\linewidth}{|p{3.5cm}|X|X|}
\hline
\rowcolor{soft}
\textbf{Type de cause} & \textbf{Description} & \textbf{Que faire ?} \\ \hline
\textbf{Cause commune} & Variation \textbf{normale}, inhérente au procédé. & Ne rien changer : le procédé est \textbf{sous contrôle}. \\ \hline
\textbf{Cause spéciale} & Variation \textbf{anormale}, due à un événement précis (panne, erreur, dérive). & \textbf{Investiguer} et prendre une action corrective. \\ \hline
\end{tabularx}
\end{center}

\vspace{0.2cm}

\begin{labo}[title={\faFlask\quad Image simple}]
\textbf{Cause commune} = la petite variation naturelle de tes mesures entre deux séries.\\
\textbf{Cause spéciale} = le jour où ton automate a été recalibré, ou quand un réactif a changé de lot.
\end{labo}

\section*{\faExclamationTriangle\quad Les règles d'alerte (Nelson / Western Electric)}

Une carte se lit avec des \textbf{règles} qui détectent des schémas anormaux. Les 8 règles de Nelson (version améliorée des règles Western Electric) :

\begin{center}
\renewcommand{\arraystretch}{1.4}
\begin{tabularx}{\linewidth}{|c|X|}
\hline
\rowcolor{soft} \textbf{N°} & \textbf{Signal d'alerte} \\ \hline
\textbf{1} & 1 point au-delà de la zone A (\(\pm 3\sigma\)). \\ \hline
\textbf{2} & 2 points sur 3 consécutifs dans la zone A ou au-delà. \\ \hline
\textbf{3} & 4 points sur 5 consécutifs dans la zone B ou au-delà. \\ \hline
\textbf{4} & 9 points consécutifs du même côté de la ligne centrale. \\ \hline
\textbf{5} & 6 points consécutifs en augmentation ou en diminution. \\ \hline
\textbf{6} & 14 points consécutifs en alternance haut-bas. \\ \hline
\textbf{7} & 15 points consécutifs dans la zone C (près de la ligne centrale). \\ \hline
\textbf{8} & 8 points consécutifs de part et d'autre de la ligne centrale, aucun en zone C. \\ \hline
\end{tabularx}
\end{center}

\vspace{0.2cm}

\begin{attention}[title={\faExclamationTriangle\quad À retenir absolument}]
La \textbf{règle n°1} est la plus importante : \textbf{un seul point hors des limites} \(\pm 3\sigma\) suffit à déclarer le procédé « hors contrôle ». Les autres règles détectent des \textbf{dérives plus subtiles}.
\end{attention}

\section*{\faChartLine\quad Les principaux types de cartes}

\begin{center}
\renewcommand{\arraystretch}{1.4}
\begin{tabularx}{\linewidth}{|p{2.5cm}|X|p{4cm}|}
\hline
\rowcolor{soft}
\textbf{Type} & \textbf{Principe} & \textbf{Quand l'utiliser ?} \\ \hline
\textbf{Shewhart} & Chaque point = une mesure ou une moyenne. Limites \(\pm 3\sigma\). & Détection de \textbf{gros déréglages} (changements brusques). \\ \hline
\textbf{EWMA} & Moyenne mobile pondérée exponentiellement. & Détection de \textbf{petites dérives} progressives. \\ \hline
\textbf{CUSUM} & Somme cumulée des écarts à la cible. & Détection de \textbf{petites dérives} persistantes. \\ \hline
\end{tabularx}
\end{center}

\begin{labo}[title={\faFlask\quad En pratique au laboratoire}]
Les cartes de \textbf{Shewhart} (\(\bar{X}\) et \(R\), ou \(\bar{X}\) et \(s\)) sont les plus utilisées en CIQ. Les cartes \textbf{EWMA} et \textbf{CUSUM} sont réservées aux cas où l'on doit détecter des dérives très faibles (ex. : méthode très stable, ou produit coûteux).
\end{labo}

\section*{\faBuilding\quad Où on l'utilise en laboratoire ?}

\begin{center}
\renewcommand{\arraystretch}{1.4}
\begin{tabularx}{\linewidth}{|p{1.2cm}|X|}
\hline
\rowcolor{soft} \textbf{\faCheck} & \textbf{Utilisation} \\ \hline
\faCheckCircle & \textbf{Contrôle qualité interne (CIQ)} : suivre les résultats d'un matériau de contrôle au quotidien. \\ \hline
\faCheckCircle & \textbf{Validation de méthode} : vérifier la stabilité sur plusieurs séries. \\ \hline
\faCheckCircle & \textbf{Comparaison inter-laboratoires} : détecter les labos déviants. \\ \hline
\faCheckCircle & \textbf{Accréditation ISO 15189 / ISO 17025} : preuve de maîtrise statistique. \\ \hline
\faCheckCircle & \textbf{Surveillance des automates} : détecter une dérive de calibration. \\ \hline
\end{tabularx}
\end{center}

\begin{retenir}[title={\faGraduationCap\quad Ce que tu dois savoir pour l'examen}]
\begin{enumerate}[leftmargin=1.2cm, itemsep=2pt]
    \item \textbf{Définir} une carte de contrôle et ses 3 lignes (LC, LSC, LIC).
    \item \textbf{Distinguer} cause commune et cause spéciale.
    \item \textbf{Citer} la règle n°1 de Nelson (point hors des limites \(\pm 3\sigma\)).
    \item \textbf{Nommer} les principaux types : Shewhart, EWMA, CUSUM.
    \item \textbf{Expliquer} l'utilité des cartes en CIQ au laboratoire.
    \item \textbf{Citer} les normes associées : ISO 7870, ISO 15189, ISO/IEC 17025.
\end{enumerate}
\end{retenir}

\vspace{0.4cm}

\begin{center}
\textit{Fiche de synthèse — Carte de contrôle — BTS BioALC — À garder dans ton classeur !}
\end{center}

\end{document}